\begin{center}
    \section*{\fontsize{20}{20}\selectfont Chapter 6}
\end{center}
\vspace{10mm}

\section{Conclusion}
The development of this AI-assisted telemedicine platform has successfully demonstrated how modern technology can bridge critical gaps in remote healthcare delivery. By integrating an intelligent symptom triage system, we created a tool that can effectively categorize and prioritize patient needs before they even enter a consultation. The seamless speech-to-text transcription feature during live sessions significantly reduces the administrative burden on doctors, allowing them to focus entirely on patient care rather than manual note-taking. Furthermore, implementing a blockchain-backed verification system for prescriptions introduces a robust layer of security and authenticity, effectively mitigating the risks associated with forged or tampered medical documents. Overall, this project proves that combining artificial intelligence with decentralized technologies can establish a more efficient, secure, and accessible healthcare experience for both patients and medical professionals.

\subsection{Limitations}
While the platform achieves its primary objectives, several limitations must be acknowledged. First, the AI symptom triage model is inherently tied to the quality and diversity of its training data; it may occasionally struggle with rare medical conditions or highly ambiguous symptoms that fall outside its established patterns. Second, the speech-to-text transcription service relies heavily on clear audio input. Background noise, overlapping speech, or poor network connectivity during online consultations can lead to transcription errors, requiring the physician to manually correct the generated reports. Additionally, the blockchain implementation for prescription verification currently operates in a controlled test environment. Deploying this on a public mainnet would require addressing potential scalability bottlenecks and ongoing transaction costs. Finally, the platform assumes a baseline level of digital literacy, which could present a barrier for elderly patients or individuals unfamiliar with smart applications.

\subsection{Future Works and Direction}
Moving forward, there are several key areas where this system can be expanded and refined. A primary focus will be on training the AI models with a broader, more diverse medical dataset to improve the accuracy of both the triage system and the automated consultation formatting. Integrating the platform with existing hospital Electronic Health Record (EHR) systems would also be a significant step, allowing doctors to access a patient's complete medical history in one place.

To improve accessibility, incorporating multi-language support for the real-time transcription and user interface will be crucial for serving a wider demographic. On the infrastructure side, we aim to transition the blockchain prescription verification from a test network to a scalable, cost-effective layer-2 solution to handle high transaction volumes without prohibitive fees. Ultimately, conducting extensive field testing and pilot programs with actual healthcare clinics will be essential. Gathering direct feedback from practicing doctors and real patients will ensure the platform not only meets technical expectations but truly enhances the daily workflow of modern medical practice.